\subsection{整环}

定义2. 环A称为整环（或整区），如果它是交换的、非零的，并且A中两个非零元素的乘积是非零的。

有理整数环Z是整环；它是交换的且非零的；两个整数 \textgreater0 的乘积是非零的；每个非零整数都具有a或-a的形式，其中a \textgreater{} 0 且 \((-a)b = -ab\) , \((-a)(-b) = ab\) 对于 a \textgreater{} 0，b \textgreater{} 0，由此得出我们的断言。

每个交换域都是整环。整环的子环是整环。特别地，交换域的子环是整环。我们将证明反之每个整环 A 都同构于某个交换域的子环。回顾（ \(\S 8\) ，第 12 号）A 已被等同于其全分式环的一个子环。于是我们的断言由以下命题得出：

命题2. 若A是整环，则A的全分式环K是交换环。

环 K 是交换的。它是非零的，因为 \(A \neq \{0\}\) 。由于 A 是整环，A 的每个非零元素都是可消去的，且 K 由形如 a/b 的分数组成，其中 \(b \neq 0\) 。现在 \(a/b \neq 0\) 蕴含 \(a \neq 0\) ，且分数 b/a 即为 a/b 的倒数。

整环的全分式环称为其分式域。这样的环与其在分式域中的像等同。

命题3. 设B为非零环，A为B的一个交换子环，使得A中每个非零元素在B中都可逆。

\begin{enumerate}
\def\labelenumi{(\alph{enumi})}
\item
  A 是一个整环。
\item
  设 \(\mathbf{A}'\) 为 \(\mathbf{A}\) 的分式域。从 \(\mathbf{A}\) 到 \(\mathbf{B}\) 的典范嵌入可以唯一地延拓为一个同构 \(f\) ，它把 \(\mathbf{A}'\) 映到 \(\mathbf{B}\)
\item
  的一个子域上。 \(f(\mathbf{A}')\) 的元素就是形如 \(xy^{-1}\) 的元素，其中 \(x \in \mathbf{A}, y \in \mathbf{A}, y \neq 0\) .
\end{enumerate}

。断言 (a) 是显然的。从 A 到 B 的典范嵌入唯一地延拓为一个同态 \(f\) ，它把 \(\mathbf{A}'\) 映入 \(\mathbf{B}\) （§ 8，第 12 号，定理 5）。断言 (b) 随后由第 1 号定理 2 得出。\(\mathbf{A}'\) 的元素是形如 \(x / y\) 的分数，其中 \(x \in \mathbf{A}, y \in \mathbf{A}, y \neq 0\) ，且 \(f(x / y) = xy^{-1}\) ，由此得 (c)。

